\section{Results}

\subsection{Finding 1: Existing Methods Are Nontrivial but Incomplete}

\begin{table*}[t]
\centering
\scriptsize
\caption{Existing/custom-stress method capability matrix. Scores come from the E47 custom stress artifacts and released-checkpoint custom-stress runs where applicable; they are not original-paper benchmark reproductions. Empty entries indicate not applicable to the method's available signal.}
\label{tab:existing-methods}
\begin{tabularx}{\textwidth}{>{\raggedright\arraybackslash}X >{\raggedright\arraybackslash}X c c c c c}
\toprule
Method & Scope & Surface inv. & Effect sens. & Auth. sens. & Resource aware & UPA \\
\midrule
Tool-name proxy & local rule baseline & 1.000 & 0.000 & 0.000 & 0.000 & 0.292 \\
Static/plan/trajectory text proxy & local text-rule baselines & 0.984 & 0.000 & 0.903 & 1.000 & 0.080 \\
Local Qwen self-audit & local LLM self-audit & 0.865 & 0.667 & 0.565 & 0.292 & 0.004 \\
TS-Guard checkpoint & released checkpoint on E47 custom stress & 0.883 & 0.625 & 0.764 & 0.833 & 0.159 \\
Safiron checkpoint & released checkpoint on E47 custom stress & 0.669 & 1.000 & 0.088 & 0.167 & 0.545 \\
Local Qwen effect/resource mapper & local LLM mapper & 0.992 & 0.833 & 0.375 & 0.125 & 0.333 \\
IPIGuard topology-only & structural topology signal & 1.000 & 0.000 & -- & -- & -- \\
CaMeL structural-policy stress & local component stress & 1.000 & -- & -- & -- & 0.250 \\
\bottomrule
\end{tabularx}
\end{table*}


Table~\ref{tab:existing-methods} shows that existing/custom-stress methods are not merely tool-name classifiers. TS-Guard reaches 0.883 surface invariance and 0.764 authorization sensitivity on the E47 custom stress, and Safiron reaches 1.000 effect sensitivity. Local Qwen variants also show nontrivial effect sensitivity. However, high performance on one axis does not imply joint binding. Safiron has 0.545 unsafe pre-allow in this custom stress; local Qwen effect/resource mapping has high surface invariance and effect sensitivity but 0.125 resource awareness and 0.333 unsafe pre-allow. These results support a diagnostic claim: current monitors can capture parts of the binding problem, but they remain incomplete when resource, authorization, and provenance must be bound jointly.

\subsection{Finding 2: The Reference Hard Guard Improves but Remains Bounded}

\begin{table}[t]
\centering
\small
\caption{Reference hard Effect-Binding Guard results. The guard improves controlled stress tradeoffs, but retains nonzero unsafe pre-allow and false denial.}
\label{tab:reference-guard}
\begin{tabular}{lccc}
\toprule
Scope & UPA & FDeny & Coverage \\
\midrule
E48 full, 822 rows & 0.036 & 0.065 & 0.917 \\
E48 held-out split & 0.051 & 0.048 & 0.955 \\
E50 source-balanced mean & 0.051 & 0.059 & 0.923 \\
\bottomrule
\end{tabular}
\end{table}


Table~\ref{tab:reference-guard} reports the reference hard guard. On the 822-row E48 suite, the guard obtains 0.917 coverage with 0.036 unsafe pre-allow and 0.065 safe false-denial. On the held-out split, unsafe pre-allow rises to 0.051 while coverage remains 0.955. Source-balanced E50 repeats preserve the same qualitative result: 0.051 mean unsafe pre-allow, 0.059 mean safe false-denial, and 0.923 mean coverage. This is evidence that explicit tuple binding is useful, but it is not a safety certificate. The guard still makes unsafe pre-commit allows and safe false denials under controlled stress.

\subsection{Finding 3: Resource/Authorization Binding Is the Main Bottleneck}

\begin{table}[t]
\centering
\small
\caption{Bottleneck stress results for the hard guard. Resource/authorization stress exposes high unsafe pre-allow despite high coverage; provenance overlay removes unsafe pre-allow in the controlled provenance stress but increases false denial and abstention.}
\label{tab:bottleneck}
\begin{tabular}{lcccc}
\toprule
Stress set & Rows & UPA & FDeny & Coverage \\
\midrule
Resource/auth stress & 240 & 0.383 & 0.000 & 0.921 \\
Provenance stress & 336 & 0.000 & 0.188 & 0.545 \\
\bottomrule
\end{tabular}
\end{table}


Table~\ref{tab:bottleneck} identifies the main bottleneck. On E50 resource/authorization stress, the hard guard has 0.383 unsafe pre-allow at 0.921 coverage. The 46 unsafe allows comprise 20 near-alias cases, 20 out-of-scope-resource cases, and six alternate-effect cases. This rules out a benign explanation in which the guard simply abstains on hard cases. Instead, it often makes a decision while missing the authorization-relevant resource relation. Provenance stress shows a different tradeoff: unsafe pre-allow is 0.000, but coverage falls to 0.545 and safe false-denial rises to 0.188. Provenance overlays can reduce unsafe pre-allow under controlled assumptions, but their utility cost must be reported rather than hidden.

\subsection{Finding 4: Authorization Infrastructure Enables Local Pre-Commit Mediation}

\begin{table}[t]
\centering
\small
\caption{Authorization-aware local pre-commit prototype. E55-v2 is used as the main result; original E55 and human-corrected sensitivity are retained as audit/sensitivity evidence.}
\label{tab:precommit}
\begin{tabular}{lcccc}
\toprule
Method / source & UPA & FDeny & Coverage & Abstain \\
\midrule
E55-v2 existing hard guard & 0.043 & 0.000 & 0.200 & 0.800 \\
E55-v2 authz-aware guard & 0.000 & 0.000 & 0.880 & 0.120 \\
Original E55 authz-aware & 0.000 & 0.000 & 0.920 & 0.080 \\
Original E55 human-corrected sensitivity & 0.000 & 0.017 & 0.920 & 0.080 \\
\bottomrule
\end{tabular}
\end{table}


Table~\ref{tab:precommit} shows the E55-v2 local pre-commit result. The existing hard guard on the same 600-row local contract has 0.200 coverage, 0.800 abstention, and 12/276 = 0.043 unsafe pre-allow. The authorization-aware prototype reaches 528/600 = 0.880 coverage and 72/600 = 0.120 abstention, while reducing unsafe pre-allow to 0/276 and safe false-denial to 0/252 under E55-v2 strict labels. The original E55 strict contract has 0.920 coverage and zero unsafe pre-allow; human-corrected sensitivity preserves zero unsafe pre-allow but introduces 4/230 = 0.017 safe false-denial. The conservative interpretation is that explicit authorization infrastructure can make local mediation more decidable, while human audit reveals residual label/atom boundary sensitivity.

\subsection{Finding 5: Atom Interfaces Carry the Improvement}

\begin{table}[t]
\centering
\small
\caption{Atom/authorization ablations from the original E55 strict result. These ablations diagnose which semantic interfaces carry the E55 improvement; the E55-v2 main table is reported separately in Table~\ref{tab:precommit}.}
\label{tab:ablation}
\begin{tabular}{lcc}
\toprule
Ablation & Main failure metric & Value \\
\midrule
No multi-resource expansion & UPA & 0.333 \\
No operation-mode binding & UPA & 0.148 \\
No provenance overlay & UPA & 0.148 \\
No alias resolution & FDeny & 0.211 \\
\bottomrule
\end{tabular}
\end{table}


Table~\ref{tab:ablation} reports original E55 strict ablations. Removing multi-resource expansion increases unsafe pre-allow to 0.333. Removing operation-mode binding or provenance overlay increases unsafe pre-allow to 0.148. Removing alias resolution does not create unsafe pre-allow in this ablation but raises safe false-denial to 0.211. These ablations support the atom-level abstraction: the improvement is not only a stricter deny rule. It depends on expanding calls into multiple resource-bearing atoms, distinguishing draft and commit modes, preserving provenance/control-source information, and resolving aliases when the contract supports it.

\subsection{Finding 6: Audits Reduce Specific Artifact Concerns}

\begin{table*}[t]
\centering
\scriptsize
\caption{Audit and validity checks. Deterministic E57-v2 checks and human audit checks answer different validity questions and are not merged into a single agreement claim.}
\label{tab:audit}
\begin{tabularx}{\textwidth}{>{\raggedright\arraybackslash}p{0.18\textwidth} >{\raggedright\arraybackslash}X >{\raggedright\arraybackslash}p{0.25\textwidth}}
\toprule
Check & Result & Scope \\
\midrule
E55-v2 leakage audit & leakage-free & deployable-input leakage diagnostic \\
E55 decision-path audit & passed & implementation-path audit \\
E56 strict label-hidden replay & passed & label-hidden replay and build audit \\
E57-v2 resource perturbation & UPA delta 0.000; coverage delta 0.000; changed decisions 0 & naming-template robustness diagnostic \\
E57-v2 reference authorizer & decision, atom-count, and atom resource-set agreement 1.000 & independent local implementation agreement \\
E57-v2 spot-audit packet & 60 rows; required fields present & human-inspectable packet \\
E57 human audit & decision agreement 54/60; all checks true 46/60 & post-hoc human audit with corrections \\
E57 human corrections & 6 decision, 13 atom, 10 violation-reason corrections & label/atom sensitivity evidence \\
\bottomrule
\end{tabularx}
\end{table*}


Table~\ref{tab:audit} separates deterministic validity checks from human audit evidence. E55-v2 leakage checks are clean, E55 decision-path audit passed, and E56 strict label-hidden replay passed. E57-v2 resource-name perturbation yields zero changed decisions and zero UPA/coverage delta for the authorization-aware guard, reducing the concern that the result is only a resource-name template artifact. The independent reference authorizer has 1.000 agreement on decisions, atom counts, and atom resource sets across the controlled E55-v2 dataset. The human audit is intentionally reported separately: it agrees on 54/60 decisions and has all checks true on 46/60 rows, while identifying six decision corrections, thirteen atom corrections, and ten violation-reason corrections. This supports corrected-label sensitivity analysis, not perfect construction agreement.
